\documentclass[12pt]{article}

% Idioma y codificación
\usepackage[spanish, es-tabla]{babel}       %es-tabla para que se titule "Tabla"
\usepackage[utf8]{inputenc}

% Márgenes
\usepackage[a4paper,top=3cm,bottom=2.5cm,left=3cm,right=3cm]{geometry}

% Comentarios de bloque
\usepackage{verbatim}

% Paquetes de links
\usepackage[hidelinks]{hyperref}    % Permite enlaces
\usepackage{url}                    % redirecciona a la web

% Más opciones para enumeraciones
\usepackage{enumitem}

% Personalizar la portada
\usepackage{titling}


% Paquetes de tablas
\usepackage{multirow}


%------------------------------------------------------------------------

%Paquetes de figuras
\usepackage{caption}
\usepackage{subcaption} % Figuras al lado de otras
\usepackage{float}      % Poner figuras en el sitio indicado H.


% Paquetes de imágenes
\usepackage{graphicx}       % Paquete para añadir imágenes
\usepackage{transparent}    % Para manejar la opacidad de las figuras

% Paquete para usar colores
\usepackage[dvipsnames]{xcolor}
\usepackage{pagecolor}      % Para cambiar el color de la página

% Habilita tamaños de fuente mayores
\usepackage{fix-cm}


%------------------------------------------------------------------------

% Paquetes de matemáticas
\usepackage{mathtools, amsfonts, amssymb, mathrsfs}
\usepackage[makeroom]{cancel}     % Simplificar tachando
\usepackage{polynom}    % Divisiones y Ruffini
\usepackage{units} % Para poner fracciones diagonales con \nicefrac

\usepackage{pgfplots}   %Representar funciones
\pgfplotsset{compat=1.18}  % Versión 1.18

\usepackage{tikz-cd}    % Para usar diagramas de composiciones
\usetikzlibrary{calc}   % Para usar cálculo de coordenadas en tikz

%Definición de teoremas, etc.
\usepackage{amsthm}
%\swapnumbers   % Intercambia la posición del texto y de la numeración

\theoremstyle{plain}

\makeatletter
\@ifclassloaded{article}{
  \newtheorem{teo}{Teorema}[section]
}{
  \newtheorem{teo}{Teorema}[chapter]  % Se resetea en cada chapter
}
\makeatother

\newtheorem{coro}{Corolario}[teo]           % Se resetea en cada teorema
\newtheorem{prop}[teo]{Proposición}         % Usa el mismo contador que teorema
\newtheorem{lema}[teo]{Lema}                % Usa el mismo contador que teorema

\theoremstyle{remark}
\newtheorem*{observacion}{Observación}

\theoremstyle{definition}

\makeatletter
\@ifclassloaded{article}{
  \newtheorem{definicion}{Definición} [section]     % Se resetea en cada chapter
}{
  \newtheorem{definicion}{Definición} [chapter]     % Se resetea en cada chapter
}
\makeatother

\newtheorem*{notacion}{Notación}
\newtheorem*{ejemplo}{Ejemplo}
\newtheorem*{ejercicio*}{Ejercicio}             % No numerado
\newtheorem{ejercicio}{Ejercicio} [section]     % Se resetea en cada section


% Modificar el formato de la numeración del teorema "ejercicio"
\renewcommand{\theejercicio}{%
  \ifnum\value{section}=0 % Si no se ha iniciado ninguna sección
    \arabic{ejercicio}% Solo mostrar el número de ejercicio
  \else
    \thesection.\arabic{ejercicio}% Mostrar número de sección y número de ejercicio
  \fi
}


% \renewcommand\qedsymbol{$\blacksquare$}         % Cambiar símbolo QED
%------------------------------------------------------------------------

% Paquetes para encabezados
\usepackage{fancyhdr}
\pagestyle{fancy}
\fancyhf{}

\newcommand{\helv}{ % Modificación tamaño de letra
\fontfamily{}\fontsize{12}{12}\selectfont}
\setlength{\headheight}{15pt} % Amplía el tamaño del índice


%\usepackage{lastpage}   % Referenciar última pag   \pageref{LastPage}
\fancyfoot[C]{\thepage}

%------------------------------------------------------------------------

% Conseguir que no ponga "Capítulo 1". Sino solo "1."
\makeatletter
\@ifclassloaded{book}{
  \renewcommand{\chaptermark}[1]{\markboth{\thechapter.\ #1}{}} % En el encabezado
    
  \renewcommand{\@makechapterhead}[1]{%
  \vspace*{50\p@}%
  {\parindent \z@ \raggedright \normalfont
    \ifnum \c@secnumdepth >\m@ne
      \huge\bfseries \thechapter.\hspace{1em}\ignorespaces
    \fi
    \interlinepenalty\@M
    \Huge \bfseries #1\par\nobreak
    \vskip 40\p@
  }}
}
\makeatother

%------------------------------------------------------------------------
% Paquetes de cógido
\usepackage{minted}
\renewcommand\listingscaption{Código fuente}

\usepackage{fancyvrb}
% Personaliza el tamaño de los números de línea
\renewcommand{\theFancyVerbLine}{\small\arabic{FancyVerbLine}}

% Estilo para C++
\newminted{cpp}{
    frame=lines,
    framesep=2mm,
    baselinestretch=1.2,
    linenos,
    escapeinside=||
}



\usepackage{listings} % Para incluir código desde un archivo

\renewcommand\lstlistingname{Código Fuente}
\renewcommand\lstlistlistingname{Índice de Códigos Fuente}

% Definir colores
\definecolor{vscodepurple}{rgb}{0.5,0,0.5}
\definecolor{vscodeblue}{rgb}{0,0,0.8}
\definecolor{vscodegreen}{rgb}{0,0.5,0}
\definecolor{vscodegray}{rgb}{0.5,0.5,0.5}
\definecolor{vscodebackground}{rgb}{0.97,0.97,0.97}
\definecolor{vscodelightgray}{rgb}{0.9,0.9,0.9}

% Configuración para el estilo de C similar a VSCode
\lstdefinestyle{vscode_C}{
  backgroundcolor=\color{vscodebackground},
  commentstyle=\color{vscodegreen},
  keywordstyle=\color{vscodeblue},
  numberstyle=\tiny\color{vscodegray},
  stringstyle=\color{vscodepurple},
  basicstyle=\scriptsize\ttfamily,
  breakatwhitespace=false,
  breaklines=true,
  captionpos=b,
  keepspaces=true,
  numbers=left,
  numbersep=5pt,
  showspaces=false,
  showstringspaces=false,
  showtabs=false,
  tabsize=2,
  frame=tb,
  framerule=0pt,
  aboveskip=10pt,
  belowskip=10pt,
  xleftmargin=10pt,
  xrightmargin=10pt,
  framexleftmargin=10pt,
  framexrightmargin=10pt,
  framesep=0pt,
  rulecolor=\color{vscodelightgray},
  backgroundcolor=\color{vscodebackground},
}

%------------------------------------------------------------------------

% Comandos definidos
\newcommand{\bb}[1]{\mathbb{#1}}
\newcommand{\cc}[1]{\mathcal{#1}}

% I prefer the slanted \leq
\let\oldleq\leq % save them in case they're every wanted
\let\oldgeq\geq
\renewcommand{\leq}{\leqslant}
\renewcommand{\geq}{\geqslant}

% Si y solo si
\newcommand{\sii}{\iff}

% Letras griegas
\newcommand{\eps}{\epsilon}
\newcommand{\veps}{\varepsilon}
\newcommand{\lm}{\lambda}

\newcommand{\ol}{\overline}
\newcommand{\ul}{\underline}
\newcommand{\wt}{\widetilde}
\newcommand{\wh}{\widehat}

\let\oldvec\vec
\renewcommand{\vec}{\overrightarrow}

% Derivadas parciales
\newcommand{\del}[2]{\frac{\partial #1}{\partial #2}}
\newcommand{\Del}[3]{\frac{\partial^{#1} #2}{\partial^{#1} #3}}
\newcommand{\deld}[2]{\dfrac{\partial #1}{\partial #2}}
\newcommand{\Deld}[3]{\dfrac{\partial^{#1} #2}{\partial^{#1} #3}}


\newcommand{\AstIg}{\stackrel{(\ast)}{=}}
\newcommand{\Hop}{\stackrel{L'H\hat{o}pital}{=}}

\newcommand{\red}[1]{{\color{red}#1}} % Para integrales, destacar los cambios.

% Método de integración
\newcommand{\MetInt}[2]{
    \left[\begin{array}{c}
        #1 \\ #2
    \end{array}\right]
}

% Declarar aplicaciones
% 1. Nombre aplicación
% 2. Dominio
% 3. Codominio
% 4. Variable
% 5. Imagen de la variable
\newcommand{\Func}[5]{
    \begin{equation*}
        \begin{array}{rrll}
            #1:& #2 & \longrightarrow & #3\\
               & #4 & \longmapsto & #5
        \end{array}
    \end{equation*}
}

%------------------------------------------------------------------------
\usepackage{extarrows} 
\usepackage{stackrel}

\begin{document}
    % 1. Foto de fondo
    % 2. Título
    % 3. Encabezado Izquierdo
    % 4. Color de fondo
    % 5. Coord x del titulo
    % 6. Coord y del titulo
    % 7. Fecha

    
    % 1. Foto de fondo
% 2. Título
% 3. Encabezado Izquierdo
% 4. Color de fondo
% 5. Coord x del titulo
% 6. Coord y del titulo
% 7. Fecha

\newcommand{\portada}[7]{

    \portadaBase{#1}{#2}{#3}{#4}{#5}{#6}{#7}
    \portadaBook{#1}{#2}{#3}{#4}{#5}{#6}{#7}
}

\newcommand{\portadaExamen}[7]{

    \portadaBase{#1}{#2}{#3}{#4}{#5}{#6}{#7}
    \portadaArticle{#1}{#2}{#3}{#4}{#5}{#6}{#7}
}




\newcommand{\portadaBase}[7]{

    % Tiene la portada principal y la licencia Creative Commons
    
    % 1. Foto de fondo
    % 2. Título
    % 3. Encabezado Izquierdo
    % 4. Color de fondo
    % 5. Coord x del titulo
    % 6. Coord y del titulo
    % 7. Fecha
    
    
    \thispagestyle{empty}               % Sin encabezado ni pie de página
    \newgeometry{margin=0cm}        % Márgenes nulos para la primera página
    
    
    % Encabezado
    \fancyhead[L]{\helv #3}
    \fancyhead[R]{\helv \nouppercase{\leftmark}}
    
    
    \pagecolor{#4}        % Color de fondo para la portada
    
    \begin{figure}[p]
        \centering
        \transparent{0.3}           % Opacidad del 30% para la imagen
        
        \includegraphics[width=\paperwidth, keepaspectratio]{assets/#1}
    
        \begin{tikzpicture}[remember picture, overlay]
            \node[anchor=north west, text=white, opacity=1, font=\fontsize{60}{90}\selectfont\bfseries\sffamily, align=left] at (#5, #6) {#2};
            
            \node[anchor=south east, text=white, opacity=1, font=\fontsize{12}{18}\selectfont\sffamily, align=right] at (9.7, 3) {\textbf{\href{https://losdeldgiim.github.io/}{Los Del DGIIM}}};
            
            \node[anchor=south east, text=white, opacity=1, font=\fontsize{12}{15}\selectfont\sffamily, align=right] at (9.7, 1.8) {Doble Grado en Ingeniería Informática y Matemáticas\\Universidad de Granada};
        \end{tikzpicture}
    \end{figure}
    
    
    \restoregeometry        % Restaurar márgenes normales para las páginas subsiguientes
    \pagecolor{white}       % Restaurar el color de página
    
    
    \newpage
    \thispagestyle{empty}               % Sin encabezado ni pie de página
    \begin{tikzpicture}[remember picture, overlay]
        \node[anchor=south west, inner sep=3cm] at (current page.south west) {
            \begin{minipage}{0.5\paperwidth}
                \href{https://creativecommons.org/licenses/by-nc-nd/4.0/}{
                    \includegraphics[height=2cm]{assets/Licencia.png}
                }\vspace{1cm}\\
                Esta obra está bajo una
                \href{https://creativecommons.org/licenses/by-nc-nd/4.0/}{
                    Licencia Creative Commons Atribución-NoComercial-SinDerivadas 4.0 Internacional (CC BY-NC-ND 4.0).
                }\\
    
                Eres libre de compartir y redistribuir el contenido de esta obra en cualquier medio o formato, siempre y cuando des el crédito adecuado a los autores originales y no persigas fines comerciales. 
            \end{minipage}
        };
    \end{tikzpicture}
    
    
    
    % 1. Foto de fondo
    % 2. Título
    % 3. Encabezado Izquierdo
    % 4. Color de fondo
    % 5. Coord x del titulo
    % 6. Coord y del titulo
    % 7. Fecha


}


\newcommand{\portadaBook}[7]{

    % 1. Foto de fondo
    % 2. Título
    % 3. Encabezado Izquierdo
    % 4. Color de fondo
    % 5. Coord x del titulo
    % 6. Coord y del titulo
    % 7. Fecha

    % Personaliza el formato del título
    \pretitle{\begin{center}\bfseries\fontsize{42}{56}\selectfont}
    \posttitle{\par\end{center}\vspace{2em}}
    
    % Personaliza el formato del autor
    \preauthor{\begin{center}\Large}
    \postauthor{\par\end{center}\vfill}
    
    % Personaliza el formato de la fecha
    \predate{\begin{center}\huge}
    \postdate{\par\end{center}\vspace{2em}}
    
    \title{#2}
    \author{\href{https://losdeldgiim.github.io/}{Los Del DGIIM}}
    \date{Granada, #7}
    \maketitle
    
    \tableofcontents
}




\newcommand{\portadaArticle}[7]{

    % 1. Foto de fondo
    % 2. Título
    % 3. Encabezado Izquierdo
    % 4. Color de fondo
    % 5. Coord x del titulo
    % 6. Coord y del titulo
    % 7. Fecha

    % Personaliza el formato del título
    \pretitle{\begin{center}\bfseries\fontsize{42}{56}\selectfont}
    \posttitle{\par\end{center}\vspace{2em}}
    
    % Personaliza el formato del autor
    \preauthor{\begin{center}\Large}
    \postauthor{\par\end{center}\vspace{3em}}
    
    % Personaliza el formato de la fecha
    \predate{\begin{center}\huge}
    \postdate{\par\end{center}\vspace{5em}}
    
    \title{#2}
    \author{\href{https://losdeldgiim.github.io/}{Los Del DGIIM}}
    \date{Granada, #7}
    \thispagestyle{empty}               % Sin encabezado ni pie de página
    \maketitle
    \vfill
}
    \portadaExamen{ffccA4.jpg}{Cálculo I\\Examen II}{Cálculo I. Examen II}{MidnightBlue}{-8}{28}{2023-2024}{Jesús Muñoz Velasco\\Arturo Olivares Martos}

    \begin{description}
        \item[Asignatura] Cálculo I.
        \item[Curso Académico] 2021-22.
        \item[Grado] Doble Grado en Ingeniería Informática y Matemáticas.
        \item[Grupo] Único.
        \item[Profesor] Jose Luis Gámez Ruiz.
        \item[Descripción] Convocatoria Extraordinaria.
        \item[Fecha] 11 de febrero de 2022.
        %\item[Duración] 3 horas.
    
    \end{description}
    \newpage

    \begin{ejercicio}[1 punto]
        Escribe las definiciones de:
        \begin{enumerate}
            \item Conjunto mayorado
            \item Sucesión convergente
            \item Sucesión de Cauchy
            \item Sucesión positivamente divergente
        \end{enumerate}
    \end{ejercicio}

    \begin{ejercicio}[1 punto]
        Enuncia los teoremas:
        \begin{enumerate}
            \item Teorema del Valor Intermedio
            \item Teorema (de compacidad) de Weierstrass
        \end{enumerate}
    \end{ejercicio}

    \begin{ejercicio}[2 puntos]
        Ayer, 10 de febrero, el día se inició con una temperatura en Granada exactamente igual a la de hoy, día 11: ambos días, a las 00:00h, 5 grados centígrados. Prueba que ayer existió una determinada hora antes del mediodía (a.m.), en la cual la temperatura coincidió exactamente con la temperatura de doce horas más tarde (misma hora numérica, pero p.m.).\\

        Sea $T:[0,24]\longrightarrow \bb{R}$ la función que relaciona cada hora con su temperatura correspondiente.\\
        
        Sea $g:[0,12]\longrightarrow \bb{R}$ dada por $g(t) = T(t)-T(t+12)$ la función que determina la diferencia de temperatura entre dos momentos del día de ayer, uno 12 horas después del otro.\\

        Con esta notación el ejercicio se traduce en demostrar que $\exists c \in [0,12] : T(c)=T(c+12)$ o, equivalentemente, que $\exists c \in [0,12] : g(t)=0$.\\

        Tenemos que la función $T$ es continua y la función $g$ también lo es (por ser diferencia de continuas). Además,
        \begin{gather*}
            g(0) = T(0)-T(12)\\
            g(12)=T(12)-T(24)
        \end{gather*}
        Como $T(24)=T(0)$ (condición dada por el enunciado), tenemos que:
        \begin{gather*}
            \left.
            \begin{array}{l}
                g(12)=T(12)-T(0)\\
                g(0)=T(0)-T(12)\\
            \end{array}
            \right\} \Rightarrow g(0) + g(12) = T(0) - T(12) +T(12) -T(0) = 0\\
        \end{gather*}

        Por tanto tenemos dos casos posibles:
        \begin{gather*}
            (1) \hspace{0.8cm} g(0)=g(12)=0 \Longrightarrow \text{ Se han encontrado dos soluciones, } \left\{
            \begin{array}{l}
                c_1=0\\
                c_2=12
            \end{array}
            \right.\\
            \vee \\
            (2) \hspace{0.4cm} \left.
            \begin{array}{l}
                g(0) \neq 0\\\\
                g(12) \neq 0 
            \end{array}
            \right\} \Longrightarrow \text{Para que su suma sea $0$, }
            \left\{
            \begin{array}{c}
                g(0)>0 \wedge g(12)<0\\
                \vee\\
                g(0)<0 \wedge g(12)>0\\
            \end{array}
            \right\}
        \end{gather*}

        En definitiva $g(0)g(12)<0$. Por el teorema (de los ceros) de Bolzano tenemos que $\exists c \in ]0,12[ : T(c)=T(c+12) \Longrightarrow \exists c \in ]0,12[ : g(t)=0$.\\

        Como $c_1,c_2$ también son posibles soluciones (caso (1)), tenemos finalmente que $\exists c \in [0,12] : T(c)=T(c+12)$.
        
            

    \end{ejercicio}

    \begin{ejercicio}[3 puntos]
        Estudia la convergencia de las siguientes sucesiones y calcula su límite (si existe):
        \begin{enumerate}
            \item $\left\{ \dfrac{1!+2!+3!+...+n!}{n!} \right\}$.\\

            Defino $\left\{ \dfrac{a_n}{b_n}\right\} = \left\{ \dfrac{1!+2!+3!+...+n!}{n!} \right\}$, donde
            \begin{gather*}
                a_n = 1!+2!+3!+...+n! \\
                b_n = n!
            \end{gather*}

            Como $\{b_n\} \nearrow \nearrow + \infty$, puedo aplicar Stolz.
            \begin{equation*}
                \dfrac{a_{n+1}-a_n}{b_{n+1}-b_n} = \dfrac{(n+1)!}{(n+1)! - n!} = \dfrac{(n+1)\cancel{n!}}{(n+1)\cancel{n!} - \cancel{n!}} = \dfrac{n+1}{n+1-1} = \dfrac{n+1}{n}\longrightarrow 1
            \end{equation*}

            Por tanto, por el criterio de Stolz,
            \[
                \left\{ \dfrac{1!+2!+3!+...+n!}{n!} \right\} \longrightarrow 1
            \]
            
            \item $\left\{ n(\sqrt[n]{a}-1) \right\}, \,\,$ $\,\, (a\in \bb{R}^+\,$ fijo$)$.\\

            Veamos primero $\{\sqrt[n]{a}\}= \{\sqrt[n]{a_n}\}$, con $a_n=a \,\, \forall n \in \bb{N}, \,\, a > 0$.
            \[
            \dfrac{a_{n+1}}{a_n}=\dfrac{a}{a}=1 \longrightarrow 1
            \]

            Por el criterio del cociente para sucesiones tenemos que $\{\sqrt[n]{a}\}\longrightarrow 1$.\\

            Aplicamos ahora el criterio de Euler:
            \begin{gather*}
                \left\{ (\sqrt[n]{a})^n\right\} = \{a\} \longrightarrow a=e^L \Longrightarrow L=\ln(a)
            \end{gather*}

            Por tanto, por Euler, $\left\{ n(\sqrt[n]{a}-1) \right\} \longrightarrow \ln(a)$\\
            
            \item $\left\{ \left( \dfrac{\alpha \sqrt[n]{a} + \beta \sqrt[n]{b}}{\alpha + \beta}\right)^n \right\}, \,\,$ $\,\, (a,b\in \bb{R}^+$ fijos $, \,\,$ $\,\, \alpha + \beta \neq 0)$.\\

            Por el ejercicio anterior, tenemos que: 
            \begin{gather*}
                \left\{ \dfrac{\alpha \sqrt[n]{a} + \beta \sqrt[n]{b}}{\alpha + \beta}\right\} \longrightarrow \dfrac{\alpha + \beta}{\alpha + \beta} = 1
            \end{gather*}

            Aplico Euler
            \begin{gather*}
                n\left( \dfrac{\alpha \sqrt[n]{a} + \beta \sqrt[n]{b}}{\alpha + \beta} -1\right) = n\left( \dfrac{\alpha \sqrt[n]{a} + \beta \sqrt[n]{b} - \alpha - \beta}{\alpha + \beta} \right) =\\\\
                = \dfrac{n\alpha (\sqrt[n]{a} -1) + n \beta (\sqrt[n]{b} -1) }{\alpha + \beta} \longrightarrow \dfrac{\alpha \ln(a) + \beta \ln(b)}{\alpha + \beta}
            \end{gather*}

            Por el criterio de Euler,
            \begin{gather*}
                \left\{ \left( \dfrac{\alpha \sqrt[n]{a} + \beta \sqrt[n]{b}}{\alpha + \beta}\right)^n \right\} \longrightarrow e^{\left(\frac{\alpha \ln(a) + \beta \ln(b)}{\alpha + \beta}\right)}
            \end{gather*}
            
        \end{enumerate}
    \end{ejercicio}

    \begin{ejercicio}[3 puntos]
        Estudia la convergencia de las series:
        \begin{enumerate}
            \item $\sum\limits_{n\geq 1} ( \sqrt{n+1} - \sqrt{n})^3$.\\
            Sea $\{b_n\} = \left\{ n ^{ -\frac{3}{2}} \right\} = \left\{ \dfrac{1}{n^{\frac{3}{2}}}\right\} = \left\{ \dfrac{1}{\sqrt{n^2}}\right\}$\\

            Sea $\{a_n\} = \left\{ (\sqrt{n+1} - \sqrt{n})^3\right\}$\\
            \begin{gather*}
                \dfrac{a_n}{b_n} = \sqrt{n^3} (\sqrt{n+1} - \sqrt{n})^3 = \left[ \sqrt[6]{n^3} (\sqrt{n+1} - \sqrt{n})\right]^3 = \\
                =\left[ \sqrt{n} (\sqrt{n+1} - \sqrt{n}) \left( \dfrac{\sqrt{n+1} + \sqrt{n}}{\sqrt{n+1} + \sqrt{n}}\right)\right]^3 = \left[ \dfrac{\sqrt{n}(\cancel{n}+1-\cancel{n})}{\sqrt{n+1}+ \sqrt{n}}\right]^3 = \\
                = \left( \dfrac{\sqrt{n}}{\sqrt{n+1} + \sqrt{n}}\right)^3 = \left( \dfrac{1}{\sqrt{\left( \frac{n+1}{n}\right) }+1}\right)^3 = \left( \dfrac{1}{\sqrt{\left( 1+ \frac{1}{n}\right) }+1}\right)^3 \longrightarrow \left( \dfrac{1}{2}\right)^3 = L
            \end{gather*}

            Como $L= \left( \dfrac{1}{2}\right)^3 = \dfrac{1}{8}\in \bb{R}^+$, por el criterio límite de comparación,
            \[
                \sum\limits_{n\geq 1} a_n \text{ converge} \Longleftrightarrow \sum\limits_{n\geq 1} b_n \text{ converge}
            \]

            Como $\sum\limits_{n\geq 1} b_n$ converge $\Longrightarrow \sum\limits_{n\geq 1} a_n$ converge
            
            \item $\sum\limits_{n\geq 1} \dfrac{n^2}{(3n+1)^2}$.\\

            Sea $\{a_n\}=\left\{ \dfrac{n^2}{(3n+1)^2} \right\} = \left\{ \dfrac{n^2}{9n^2 + 6n + 1}\right\} \longrightarrow \dfrac{1}{9} \neq 0$\\

            Como $\{a_n\} \nrightarrow 0 \Longrightarrow$ No cumple la condición necesaria de convergencia.\\

            Por tanto $\sum\limits_{n\geq 1} \dfrac{n^2}{(3n+1)^2}$ no converge.
            
            
            \item $\sum\limits_{n \geq 1} \dfrac{\cos^3(n^2 + 7n - 10)}{n^3}$.\\

            Defino $\sum\limits_{n\geq 1} a_n = \sum\limits_{n \geq 1} \dfrac{\cos^3(n^2 + 7n - 10)}{n^3}$.

            Como el coseno puede ser negativo, no es una serie de términos no negativos. Veamos si converge absolutamente. \\
            
            Estudiamos por tanto la convergencia de $\sum\limits_{n\geq 1} |a_n| = \sum\limits_{n \geq 1} \dfrac{|\cos^3(n^2 + 7n - 10)|}{n^3}$.\\

            Sea $\{b_n\}=\left\{ \dfrac{1}{n^3}\right\}$.
            \begin{gather*}
                |a_n|\leq b_n \Longleftrightarrow \dfrac{|\cos^3(n^2 + 7n - 10)|}{\cancel{n^3}} \leq \dfrac{1}{\cancel{n^3}} \Longleftrightarrow |\cos^3(n^2 + 7n - 10)|\leq 1 \Longleftrightarrow\\ \\
                \Longleftrightarrow |\cos(n^2 + 7n - 10)|\leq 1 \hspace{1cm} \text{ Cierto }\forall n \in \bb{N}
            \end{gather*}

            Por tanto,
            \[
                0\leq |a_n| \leq b_n \,\,\, \forall n \in \bb{N}
            \]

            Por el criterio de comparación, $\sum\limits_{n \geq 1}b_n$ converge $\Longrightarrow \sum\limits_{n\geq 1} |a_n| $ converge\\

            Por tanto, $\sum\limits_{n\geq 1} |a_n| $ converge $\Longrightarrow \sum\limits_{n\geq 1} a_n$ converge absolutamente $\Longrightarrow \sum\limits_{n\geq 1} a_n $ converge 

            
        \end{enumerate}
    \end{ejercicio}


     
\end{document}
